\section{Module \texttt{input.c}}

Le module \texttt{input.c} regroupe les composants de saisie. Il encapsule trois widgets principaux : l'entrée de texte, la liste déroulante moderne et le bouton numérique à incréments.

\subsection{Fonction \texttt{create\_entry}}

Cette fonction prépare un \texttt{GtkEntry} puis applique le placeholder, le texte initial, le type de saisie attendu et les éventuelles contraintes de visibilité ou de longueur.

\begin{lstlisting}[language=C]
GtkWidget *create_entry(const entry_config *config) {
    GtkWidget *entry = gtk_entry_new();

    if (config == NULL) {
        return entry;
    }

    if (config->placeholder) {
        gtk_entry_set_placeholder_text(GTK_ENTRY(entry), config->placeholder);
    }
    if (config->default_text) {
        gtk_editable_set_text(GTK_EDITABLE(entry), config->default_text);
    }
    gtk_entry_set_input_purpose(GTK_ENTRY(entry), parse_input_purpose(config->purpose_hint));
    gtk_entry_set_visibility(GTK_ENTRY(entry), !config->password);
    gtk_editable_set_editable(GTK_EDITABLE(entry), config->editable);
    if (config->max_length > 0) {
        gtk_entry_set_max_length(GTK_ENTRY(entry), config->max_length);
    }
    apply_widget_style(entry, &config->style);

    if (config->on_changed) {
        g_signal_connect(entry, "changed", G_CALLBACK(config->on_changed), config->user_data);
    }

    return entry;
}
\end{lstlisting}

\subsection{Fonction \texttt{create\_dropdown}}

\begin{lstlisting}[language=C]
GtkWidget *create_dropdown(const dropdown_config *config) {
    GtkWidget *dropdown;

    if (config == NULL || config->options == NULL) {
        return gtk_drop_down_new(NULL, NULL);
    }

    dropdown = gtk_drop_down_new_from_strings(config->options);
    gtk_drop_down_set_selected(GTK_DROP_DOWN(dropdown), config->selected_index);
    gtk_drop_down_set_enable_search(GTK_DROP_DOWN(dropdown), config->enable_search);
    apply_widget_style(dropdown, &config->style);

    if (config->on_selected) {
        g_signal_connect(dropdown, "notify::selected", G_CALLBACK(config->on_selected), config->user_data);
    }

    return dropdown;
}
\end{lstlisting}

\subsection{Fonction \texttt{create\_spin\_button}}

\begin{lstlisting}[language=C]
GtkWidget *create_spin_button(const spin_button_config *config) {
    GtkAdjustment *adj;
    GtkWidget *spin;

    if (config == NULL) {
        adj = gtk_adjustment_new(0.0, 0.0, 100.0, 1.0, 5.0, 0.0);
        return gtk_spin_button_new(adj, 1.0, 0);
    }

    adj = gtk_adjustment_new(config->value, config->min_value, config->max_value, config->step, config->page, 0.0);
    spin = gtk_spin_button_new(adj, config->step, config->digits);
    gtk_spin_button_set_numeric(GTK_SPIN_BUTTON(spin), config->numeric_only);
    gtk_spin_button_set_wrap(GTK_SPIN_BUTTON(spin), config->wrap);
    apply_widget_style(spin, &config->style);

    if (config->on_value_changed) {
        g_signal_connect(spin, "value-changed", G_CALLBACK(config->on_value_changed), config->user_data);
    }

    return spin;
}
\end{lstlisting}

\subsection{APIs GTK utilisées}

\begin{center}
\begin{tabular}{|l|p{8cm}|}
\hline
\textbf{API} & \textbf{Rôle} \\
\hline
\texttt{gtk\_entry\_set\_placeholder\_text()} & Définit le texte d'indication d'un champ vide. \\
\hline
\texttt{gtk\_entry\_set\_input\_purpose()} & Informe GTK du type de saisie attendu. \\
\hline
\texttt{gtk\_drop\_down\_new\_from\_strings()} & Construit une liste de choix depuis un tableau de chaînes. \\
\hline
\texttt{gtk\_drop\_down\_set\_enable\_search()} & Active la recherche intégrée du widget déroulant. \\
\hline
\texttt{gtk\_adjustment\_new()} & Définit l'intervalle numérique et les pas du contrôle. \\
\hline
\texttt{gtk\_spin\_button\_new()} & Crée le widget de saisie numérique associé à un ajustement. \\
\hline
\end{tabular}
\end{center}

\newpage
